\begin{frame}[fragile]{Three Options for Configuration}
  \begin{description}[Configuration Files]
    \item[In-Code] configuration is provided for \alert{anonymous and named} instances.
      \begin{itemize}
        \item Configure \alert{port number} of a server instance-object \alert{\lstinline|echoServer|}
        \begin{cppcode}[autogobble=true,basicstyle=\normalsize\codefontfamily, numberstyle=\normalsize\numberfontfamily,frame=none,numbers=none,numbersep=0pt]
            EchoServerType echoServer("echo");
            echoServer.getConfig().Local::setPort(8080);
        \end{cppcode}
      \end{itemize}
    \item[Commandline] configuration exists for \alert{named instances only}
      \begin{itemize}
        \item Configure the \alert{port number} of instance named \alert{\lstinline|echo|} in application \alert{echoserver}
        \begin{cppcode}[autogobble=true,basicstyle=\normalsize\codefontfamily, numberstyle=\normalsize\numberfontfamily,frame=none,numbers=none,numbersep=0pt]
            command@line~> ./echoserver echo local --port 8080
        \end{cppcode}
      \end{itemize}
    \item[Configuration File] configuration exists for \alert{named instances only}
      \begin{itemize}
        \item Configure the \alert{port number} of instance named \alert{\lstinline|echo|} in application \alert{echoserver}\\
        \medskip
        \begin{tabular}{|l|}
        \hline
        \lstinline|~/.config/snode.c/echoserver.cfg|\\
        \specialrule{.1em}{.0em}{.0em}
        \lstinline|echo.local.port=8080|\\
        \hline
        \end{tabular}
      \end{itemize}
  \end{description}
    \begin{itemize}
      \item \alert{Configuration File entries} takes precedence over \alert{In-Code} configuration and
      \item \alert{Command-Line options} takes precedence over \alert{Configuration File} configuration.
    \end{itemize}
\end{frame}
